\documentclass[12pt]{article}

%% Basic document formatting
\usepackage{amsmath, amsthm, amssymb, amsfonts}
\usepackage{mathtools}
\usepackage{xspace}
\usepackage{thmtools}
\usepackage{graphicx}
\usepackage{setspace}
\usepackage{fancyhdr}
\usepackage{titling}
\usepackage{hyperref}
\usepackage[left=0.4in,right=0.4in,top=1in,bottom=1in]{geometry}
\usepackage{float}
\usepackage{tabularx}
\usepackage[utf8]{inputenc}
\usepackage[english]{babel}
\usepackage{framed}
\usepackage[dvipsnames]{xcolor}
\usepackage{environ}
\usepackage{tcolorbox}
\tcbuselibrary{theorems,skins,breakable}

\usepackage{enumitem}
\setlist[enumerate]{leftmargin=*}
\setlist[enumerate,1]{labelindent=\parindent}
\setlist[enumerate,2]{labelindent=0pt}

% Blackboard Bold
\newcommand{\Z}{\mathbb{Z}}                 % integers
\newcommand{\Zpl}{\mathbb{Z}_{+}}           % positive integers
\newcommand{\Znn}{\mathbb{Z}_{\geq 0}}      % nonnegative integers. 
\newcommand{\Q}{\mathbb{Q}}                 % rationals
\newcommand{\Qpl}{\mathbb{Q}_{+}}           % positive Rationals
\newcommand{\R}{\mathbb{R}}                 % reals
\newcommand{\Rpl}{\mathbb{R}_{+}}           % positive Reals
\newcommand{\C}{\mathbb{C}}                 % complex numbers
\newcommand{\F}{\mathbb{F}}                 % field

% Words
\newcommand{\st}{\text{ such that }}
\newcommand{\wrt}{\text{ with respect to }}
\newcommand{\with}{\text{ with }}
\newcommand{\ie}{\text{, i.e. }}
\newcommand*{\ditto}{\text{---}\texttt{"}\text{---}}

% Operators
\newcommand{\abs}[1]{\left|#1\right|}                   % absolute value
\newcommand{\norm}[1]{\abs{\abs{#1}}}
\newcommand{\floor}[1]{\left\lfloor #1 \right\rfloor}   % floor
\newcommand{\ceil}[1]{\left\lceil #1 \right\rceil}      % ceiling
\newcommand{\powset}[1]{\mathscr{P}(#1)}

% Category Theory 
\renewcommand{\hom}{\text{Hom}}
\newcommand{\homspace}[3]{\hom_{#1}(#2, #3)}
\newcommand{\coproduct}[9]{
  \begin{tikzcd}[ampersand replacement=\&]
      \& \& #1 \arrow[dll, bend right, "#6"] \arrow[dl, "#5"] \\
   #4 \& #3 \arrow[l, "#9"] \\
      \& \& #2 \arrow[ull, bend left, "#8"] \arrow[ul, "#7"]
  \end{tikzcd}
}

\newcommand{\product}[9]{ 
  \begin{tikzcd}[ampersand replacement=\&]
      \& \& #3 \\
      #1 \arrow[r, "#7"] \arrow[urr, bend left, "#5"] \arrow[drr, bend right, "#6"] \& #2 \arrow[ur, "#8"] \arrow[dr, "#9"] \\ 
      \& \& #4
  \end{tikzcd}
}


% Algebra 
\newcommand{\Syl}[2]{\text{Syl}_{#1}{(#2)}}           % set of Sylow-p subgroups 
\newcommand{\<}{\langle}                % \<x\>, subgroup generated by x 
\renewcommand{\>}{\rangle}
\renewcommand{\index}[2]{[#1 : #2]}
\newcommand{\id}{\text{id}}             % identity element
\newcommand{\lhdneq}{%
  \mathrel{\ooalign{$\lneq$\cr\raise.22ex\hbox{$\lhd$}\cr}}}
\newcommand{\order}[1]{\text{o}(#1)}    % order of an element
\newcommand{\spec}{\operatorname{Spec}}
\newcommand{\rad}{\operatorname{rad}}   % radical of an ideal 
\newcommand{\sym}[1]{\mathfrak{S}_{#1}}
\newcommand{\aut}[1]{\text{Aut}(#1)} 
\newcommand{\gal}[2]{\text{Gal}(#1 / #2)} 
\newcommand{\inn}[1]{\text{Inn}(#1)} 
\let\oldcong\cong
\let\oldequiv\equiv
\renewcommand{\cong}{\oldequiv}
\renewcommand{\equiv}{\oldcong}

\newcommand{\triangleleftneq}{\mathrel{\ooalign{%
  $\trianglelefteq$\cr
  \hidewidth\raise0.06ex\hbox{$\not\mkern4mu$}\hidewidth\cr
}}}

\makeatletter % cycle
\newcommand{\cyc}[1]{(\mathbf{\cyc@process#1\relax})}
\def\cyc@process#1#2\relax{%
	#1%
	\ifx\relax#2\relax
	\else
		\,\cyc@process#2\relax
	\fi
}
\makeatother

\newcommand{\GL}[2]{\text{GL}_{#1}(#2)}                             % general linear group
\newcommand{\SL}[2]{\text{SL}_{#1}(#2)}                             % special linear group
\newcommand{\gl}[2]{\mathfrak{gl}_{#1}(#2)}
\renewcommand{\sl}[2]{\mathfrak{sl}_{#1}(#2)}
\newcommand{\so}[2]{\mathfrak{so}_{#1}(#2)}
\newcommand{\su}[1]{\mathfrak{su}(#1)}

% Linear Algebra
\newcommand{\bmat}[1]{\begin{bmatrix}#1\end{bmatrix}}   % bracketed matrix
\newcommand{\rank}{\operatorname{rank}}                 % rank
\newcommand{\nullity}{\operatorname{nullity}}           % nullity
\newcommand{\tr}{\operatorname{tr}}

% Combinatorics
\newcommand{\qnum}[1]{\left[#1\right]_q}                % q-analog
\newcommand{\qfact}[1]{\left[#1\right]!_q}              % q-analog factorial 
\newcommand{\qbinom}[2]{\binom{#1}{#2}_q}               % Gaussian binomial coefficient

% Topology/Analysis
\newcommand{\ball}[2]{\text{B}_{#1}(#2)}  % B_r(x): open r-balls around x
\newcommand{\diam}{\text{diam}}           % diamter of a set in metric space 
\newcommand{\lowint}[2]{\underline{\int_{#1}^{#2}}}
\newcommand{\upint}[2]{\overline{\int}_{#1}^{#2}}

% Misc Notation
\newcommand{\oneton}[1]{\{1, \ldots, #1\}}
\newcommand{\defeq}{\vcentcolon=}   % :=
\newcommand{\eqdef}{=\vcentcolon}   % =: 
\renewcommand{\bf}[1]{\textbf{#1}}
\newcommand{\im}{\operatorname{im}}
\newcommand{\fnrest}[2]{\left.#1\right|_{#2}}
\newcommand{\isom}{\overset{\sim}{\rightarrow}} 


\renewcommand{\thesection}{\arabic{section}.}
\renewcommand{\thesubsection}{(\alph{subsection})}

\newtheorem{claim}{Claim}
\newtheorem*{lemma}{Lemma}

%% Headers & title setup
\newcommand{\course}{Math100B}
\newcommand{\myname}{Jay Ser}
\setlength{\headheight}{14.5pt}
\pagestyle{fancy}
\fancyhf{}
\renewcommand{\headrulewidth}{0.4pt}
\lhead{\course}
\rhead{\myname}
\cfoot{\thepage}
\setlength{\droptitle}{-4em} 
\title{\course\ - HW \#9} 
\author{\myname}
\date{2026.03.12}

\begin{document}
\maketitle
\thispagestyle{fancy}

%------------------------------------------------------------------------------%

\section{}
Suppose there are $x_i \in F(\zeta_3 \sqrt[3]{2})$ satisfying 
\begin{equation} \label{eq:1_sln}
  x_{1}^{2} + \ldots + x_{n}^{2} = -1
\end{equation}
Let 
$$\varphi: F(\zeta_3 \sqrt[3]{2}) \rightarrow F(\sqrt[3]{2}); \: \left. \varphi \right|_\Q = \id, \, \zeta_3 \sqrt[3]{2} \mapsto \sqrt[3]{2}$$ 
be the isomorphism between the two field extensions. 
Applying $\varphi$ to \eqref{eq:1_sln}, obtain 
$$y_{1}^{2} + \ldots + y_{n}^{2} = -1,$$ 
where $y_i = \varphi(x_i) \in \Q(\sqrt[3]{2}) \subset \R$. 
This is clearly a contradiction since each $y_i \in \R$ and hence each $y_{i}^{2} \geq 0$. 

\section{} 
$\alpha^2 \in F(\alpha)$, so $F(\alpha^2) \subset F(\alpha)$ and $[F(\alpha, \alpha^2) : F(\alpha)] = 1$. 
Particularly, 
$$[F(\alpha, \alpha^2) : F] = [F(\alpha, \alpha^2) : F(\alpha)] [F(\alpha) : F] = 5$$ 
Suppose $\alpha \notin F(\alpha^2)$. 
Then $[F(\alpha, \alpha^2) : F(\alpha^2)] > 1$. 
By 
$$5 = [F(\alpha, \alpha^2) : F] = [F(\alpha, \alpha^2) : F(\alpha^2)] [F(\alpha^2) : F],$$ 
$[F(\alpha^2) : F] = 1$, i.e., $\alpha^2 \in F$.  
Then $f(x) = x^2 - \alpha^2$ is the minimal monic polynomial over $F$ satisfied by $\alpha$ (since $f$ is degree 2, and clearly no polynomial of degree 1 has $\alpha$ as a root).   
Of course, this contradicts the assumption that $[F(\alpha) : F] = 5$, or in other words, the minimal polynomial for $\alpha$ over $F$ is degree 5. 
Hence $\alpha \in F(\alpha^2)$, which shows $F(\alpha) = F(\alpha^2)$. 

\newcommand{\fieldK}{\Q(\sqrt[3]{2})}
\section{}
$f(x) \defeq x^4 + 3x + 3$ is irreducible over $\Q$ by the Eisenstein criterion. 

Suppose $f$ is reducible over $\fieldK$. 
Let $g \in \fieldK[x]$ be a nonconstant irreducible factor of $f$ such that $1 \leq \deg(g) \leq 3$. 
Denote $d \defeq \deg(g)$. 
There is a field $L$ such that $\alpha \in L$ is a root of $g$. 
So $[\Q(\sqrt[3]{2}, \alpha) : \fieldK] = d$. 
Since $g$ is a factor of $f$, $\alpha$ is a root of $f$ as well, and since $f$ is irreducible over $\Q$, $[\Q(\alpha) : \Q] = 4$. 
Finally, $[\fieldK : \Q] = 3$. 
Hence $[\Q(\sqrt[3]{2}, \alpha), \Q] = 3d$. 
Putting this all togother, 
$$3d = [\Q(\sqrt[3]{2}, \alpha) : \Q] = [\Q(\sqrt[3]{2}, \alpha) : \Q(\alpha)][\Q(\alpha): \Q],$$
i.e., $4 \mid 3d$. 
But this is impossible since $1 \leq d \leq 3$. 
This contradiction shows that $f$ is irreducible over $\fieldK$. 

\section{} 
Recall that the minimal polynomial for $\zeta_5$ is $x^4 + x^3 + x^2 + x + 1$ and that for $\zeta_7$ is $x^6 + x^5 + x^4 + x^3 + x^2 + x + 1$.
So $[\Q(\zeta_5) : \Q] = 4$ and $[\Q(\zeta_7) : \Q] = 6$. 
Suppose $\zeta_5 \in \Q(\zeta_7)$, i.e., $\Q(\zeta_5) \subset \Q(\zeta_7)$. 
Then 
$$6 = [\Q(\zeta_7) : \Q] = [\Q(\zeta_7) : \Q(\zeta_5)] [\Q(\zeta_5) : \Q] = 4 [\Q(\zeta_7) : \Q(\zeta_5)]$$ 
which contradicts $4 \nmid 6$. 

\section{}
$\zeta_4$ satisfies $x^4 - 1 = (x^2 - 1)(x^2 + 1)$. 
$\zeta_{4}^{2} = \cos(2(2\pi) / 4) + i\sin(2(2 \pi) / 4) = -1$, so $\zeta_4$ is the root of $x^2 + 1$. 
$x^2 + 1$ is irreducible over $\Q$ by the rational root test;
the minimal polynomial for $\zeta_4$ over $\Q$ is $x^2 + 1$. 
Of course, $\zeta_4$ still satisfies $x^2 + 1$ over $\Q(\zeta_3)$. 
To see that $x^2 + 1$ is the minimal polynomial for $\zeta_4$ over $\Q(\zeta_3)$, suppose it is not; 
then necessarily $[\Q(\zeta_4, \zeta_3) : \Q(\zeta_3)] = 1$. 
Since 
$$[\Q(\zeta_3, \zeta_4) : \Q] = [\Q(\zeta_3, \zeta_4) : \Q(\zeta_3)][\Q(\zeta_3) : \Q] = [\Q(\zeta_3, \zeta_4) : \Q(\zeta_4)][\Q(\zeta_4) : \Q],$$
it follows that $[\Q(\zeta_3, \zeta_4) : \Q(\zeta_4)] = 1$,
i.e., $\zeta_3 \in \Q(\zeta_4) \equiv \Q(i)$, where the isomorphism due to the fact that $x^2 + 1$ is the minimal polynomial for $\zeta_4$ and $i$ over $\Q$. 
Hence 
$$\zeta_3 = -\frac{1}{2} + i \frac{\sqrt{3}}{2} = a + bi$$
for $a, b \in \Q$, which is clearly a contradiction. 
This shows that $x^2 + 1$ is still the minimal polynomial for $\zeta_4$ over $\Q(\zeta_3)$. 

We already know that the minimal polynomial for $\zeta_5$ over $\Q$ is $x^4 + x^3 + x^2 + x + 1$ since 5 is prime. 
To determine the degree of $\zeta_5$ over $\Q(\zeta_3)$, notice that $\Q(\zeta_3, \zeta_5) = \Q(\zeta_{15})$: 
$\zeta_{15}^{3} = \zeta_5$ and $\zeta_{15}^{5} = \zeta_3$ shows $\Q(\zeta_3, \zeta_5) \subset \Q(\zeta_{15})$. 
The other inclusion follows from $(\zeta_3 \zeta_5)^2 = e^{32\pi / 15} = e^{2\pi / 15} = \zeta_{15}$. 
The word around the street is that the minimal polynomial for $\zeta_{15}$ has degree 8. 
This gives 
$$8 = [\Q(\zeta_{15}) : \Q] = [\Q(\zeta_{3}, \zeta_5) : \Q(\zeta_3)][\Q(\zeta_3) : \Q] = 2[\Q(\zeta_3, \zeta_5) : \Q(\zeta_3)],$$
which mean that the degree of $\zeta_5$ over $\Q(\zeta_3)$ is 4. 
It follows that $x^4 + x^3 + x^2 + x + 1$ is still the minimal polynomial for $\zeta_5$ over $\Q(\zeta_3)$.  

$\zeta_6$ satisfies $x^6 - 1 = (x^3 - 1)(x^3 + 1)$. 
$\zeta_{6}^{3} = \cos{\pi} + i\sin{\pi} = -1$. 
Hence $\zeta_6$ satisfies $x^3 + 1 = (x + 1)(x^2 - x + 1)$. 
Of course, $\zeta_6$ doesn't satisfy $x + 1$, hence it satisfies $(x^2 - x + 1)$, which is irreducible over $\Q$ by the rational root test. 
So the minimal polynomial for $\zeta_6$ over $\Q$ is $x^2 - x + 1$. 
Next, because $(\zeta_6)^2 = \zeta_3$, $[\Q(\zeta_3, \zeta_6) : \Q(\zeta_6)] = 1$. 
Thus the equality 
$$[\Q(\zeta_3, \zeta_6) : \Q] = [\Q(\zeta_3, \zeta_6) : \Q(\zeta_3)] [\Q(\zeta_3) : \Q] =   [\Q(\zeta_3, \zeta_6) : \Q(\zeta_6)] [\Q(\zeta_6) : \Q]$$
and that $[\Q(\zeta_6) : \Q] = [\Q(\zeta_3) : \Q] = 2$ implies $[\Q(\zeta_3, \zeta_6) : \Q(\zeta_3)] = 1$, 
i.e., the minimal polynomial for $\zeta_6$ over $\Q(\zeta_3)$ is just $x - \zeta_6$. 

\section{}
$a^{1/4}$ satisfies $x^4 - a$, so $[\Q(a^{1/4}) : \Q] \leq 4$. 
Also, $a^{1/2}$ satisfies $x^2 - a$, which is irreducible by the assumption that it has no roots, hence $[\Q(a^{1/2}) : \Q] = 2$. 
$(a^{1/4})^{2} = a^{1/2}$, so $\Q(a^{1/2}) \subset \Q(a^{1/4})$, which gives 
$$4 \geq [\Q(a^{1/4}) : \Q] = [\Q(a^{1/4}) : \Q(a^{1/2})] [\Q(a^{1/2}) : \Q] = 2 [\Q(a^{1/4}) : \Q(a^{1/2})].$$
If $[\Q(a^{1/4}) : \Q] = 2$, then $[\Q(a^{1/4}) : \Q(a^{1/2})] = 1$, implying $a^{1/4} \in \Q(a^{1/2})$.  
$\Q(a^{1/2}) \equiv \{\alpha + \beta a^{1/2} \mid \alpha, \beta \in \Q\}$, so 
$$\alpha^{1/4} = \alpha + \beta a^{1/2} \implies a^{1/2} = \alpha^2 + 2\alpha \beta \alpha^{1/2} + \beta^2 a$$ 
for some $\alpha, \beta \in \Q$, which contradicts assumption that $a^{1/2} \notin \Q$. 
It follows that $[\Q(a^{1/4}) : \Q] = 4$, as was to be shown. 

\section{}
I show ``$f$ irreducible over $L \implies g$ irreducible over $K$." 
The reverse direction follows by the exact same logic. 
Denote $n \defeq \deg(f)$ and $m \defeq \deg(g)$. 
If $f$ is irreducible over $L$, then $\Q(\alpha, \beta) \equiv L[x] / \<f(x)\>$, hence 
\begin{align*}
  nm & = [\Q(\alpha, \beta) : \Q(\beta)][\Q(\beta) : \Q] \\ 
     & = [\Q(\alpha, \beta) : \Q] \\ 
     & = [\Q(\alpha, \beta) : \Q(\alpha)][\Q(\alpha) : \Q] \\ 
     & = n [\Q(\alpha, \beta) : \Q(\alpha)]
\end{align*}
This shows that $[\Q(\alpha, \beta) : \Q(\alpha)] = m$. 
Of course, $\beta$ still satisfies $g(x) \in K[x]$ with $\deg(g) = m$.
Thus $g$ is necessarily the minimal polynomial for $\beta$ over $K$. 
It must be that $g$ is irreducible over $K$. 

\section{} 
\subsection{} 
Let $F' \defeq F(\alpha) \cap F(\beta)$. 
It is trivial to show that in general, the intersection of subfields is a subfield as well. 
$F(\alpha)$ and $F(\beta)$ are subfields of $E$, so $F'$ is a subfield of $E$. 
Namely, $F'$ is a field. 

\subsection{}
Let $n \defeq [F(\alpha) : F]$ and $m \defeq [F(\beta) : F]$. 
Since $F \subsetneq F'$ and $F'$ is a subfield of $E$, 
\begin{equation} \label{eq:prop_ext}
  t \defeq [F' : F] > 1
\end{equation}
Since $F' \subset F(\alpha)$, 
\begin{equation} \label{eq:big_tower}
  [F(\alpha, \beta) : F] = [F(\alpha, \beta) : F'] [F' : F]
\end{equation}
Also, notice $[F(\alpha) : F] = [F(\alpha) : F'] [F' : F]$ and thus $[F(\alpha) : F'] = n / t$. 
Similarly, $[F(\beta) : F'] = m / t$, which means $[F(\alpha, \beta), F'] \leq nm / t^2$. 
Combining this result with \eqref{eq:big_tower},
$$[F(\alpha, \beta) : F] \leq nm / t.$$ 
The inequality \eqref{eq:prop_ext} shows $[F(\alpha, \beta), F] < nm$, as was to be demonstrated.  

\end{document}
